{% autoescape off %}
\documentclass{article}

\usepackage{fancyhdr}

\pagestyle{fancy}
\fancyhf{}
\renewcommand{\headrulewidth}{0pt}

\usepackage{xcolor}
\usepackage{array}
\usepackage[utf8]{inputenc}
\usepackage[labelformat=empty,skip=4pt]{caption}
\usepackage[export]{adjustbox}
\usepackage{wrapfig}
\usepackage[a6paper,landscape]{geometry}
\usepackage{qrcode}
\usepackage[swedish]{isodate}
\usepackage[parfill]{parskip}
\geometry{tmargin=1cm,lmargin=1.5cm,rmargin=1.5cm}
\isodate

\def\version{0.1}

%\ifdefined\tagid
%\else
%	\def\tagid{dummy tag}
%\fi

\def\tagid{ {{ tag.pk }} }
\def\tagurl{ {{ url }} }
\def\memberid{ {{ tag.member_id.id }} }
\def\tagboxnumber{ {{ tag.box_number }} }
\def\membername{ {{ tag.member_id.name }} }
\def\tagprintdate{ {{ tag.print_date }} }
\def\tagcomment{ {{ tag.comment }} }

\lfoot{v.\version}
\rfoot{\today}

\begin{document}

\begin{wrapfigure}[0]{r}{0.3\textwidth}
\vspace{0.5cm}
\qrcode[height=4cm]{\tagurl}
\caption{\tagurl}
\end{wrapfigure}

\includegraphics[width=0.7\textwidth]{MakersLink-line-color.png}

\vspace{1.3cm}

Medlems-nr:
\begin{tabular}{>{\centering}p{3em}}
	\memberid
	\tabularnewline
    \hline

\end{tabular}Låd nr: \begin{tabular}{>{\centering}p{3em}}
	\tagboxnumber
	\tabularnewline
    \hline

\end{tabular}

Namn:
\begin{tabular}{>{\centering}p{13em}}
	\membername
	\tabularnewline
    \hline

\end{tabular}

Datum:
\begin{tabular}{>{\centering}p{12.5em}}
	\tagprintdate
	\tabularnewline
    \hline

\end{tabular}

Kommentar:
\begin{tabular}{>{\centering}p{10.5em}}
	\tagcomment
	\tabularnewline
    \hline
\end{tabular}

\begin{tabular}{>{\centering}p{16.5em}}
	\tabularnewline
    \hline
\end{tabular}
\newpage

\section*{Regler}

\begin{itemize}
	\item Alla lådor ska vara märkta med en etikett.
    \item Alla föremål som går ner i en låda ska vara i en låda.
    \item Alla föremål som inte går ner i en låda ska vara märkta med en etikett.
\end{itemize}

\section*{Användning}

Denna etikett är till för att märka lådor och föremål som medlemmar förvarar i makerspacets lokal.

\begin{enumerate}
  \item Fyll i uppgifter på sidan med logotypen.
  \item Om etiketten är utskriven på A4 vik den dubbel så att sidan med logotypen är synlig.
  \item Fäst etiketten på en kortsida på lådan antingen i  en plastficka eller med tejp.
\end{enumerate}

\end{document}
{% endautoescape %}
